\documentclass[11pt]{article}
\usepackage[margin=1in]{geometry}
\usepackage{amsmath,amssymb,amsthm,mathtools,booktabs}
\usepackage{microtype}
\usepackage[colorlinks=true,linkcolor=blue,citecolor=blue,urlcolor=blue]{hyperref}
\hypersetup{pdftitle={An explicit rate for ultraflat real Littlewood polynomials},pdfauthor={Jeffery Kline},pdfsubject={Version 0.1.0, release candidate}}
\setlength{\emergencystretch}{2em}
\allowdisplaybreaks[2]
\newtheorem{theorem}{Theorem}[section]
\newtheorem{proposition}[theorem]{Proposition}
\newtheorem{lemma}[theorem]{Lemma}
\newtheorem{corollary}[theorem]{Corollary}
\theoremstyle{remark}
\newtheorem{remark}[theorem]{Remark}
\newcommand{\R}{\mathbb R}
\newcommand{\C}{\mathbb C}
\newcommand{\Z}{\mathbb Z}
\newcommand{\T}{\mathbb T}
\newcommand{\E}{\mathbb E}
\newcommand{\Prob}{\mathbb P}
\newcommand{\ee}{\mathrm e}
\newcommand{\sgn}{\operatorname{sgn}}
\newcommand{\codeg}{\operatorname{codeg}}
\newcommand{\supp}{\operatorname{supp}}
\newcommand{\ceil}[1]{\left\lceil #1\right\rceil}
\newcommand{\floor}[1]{\left\lfloor #1\right\rfloor}
\newcommand{\abs}[1]{\left|#1\right|}
\newcommand{\norm}[1]{\left\|#1\right\|}
\newcommand{\Ethree}{\mathcal E_3}
\title{An explicit rate for ultraflat real Littlewood polynomials}
\author{Jeffery Kline}
\date{October 10, 2026\\[3pt]\small Version 0.1.0, release candidate}
\begin{document}
\maketitle
\begin{abstract}
A real Littlewood polynomial of length $N$ is a polynomial
$\sum_{j=0}^{N-1}\epsilon_jz^j$ whose coefficients are all $+1$ or $-1$.
OpenAI's preprint \emph{Ultraflat real Littlewood polynomials}
(October 5, 2026) proved that for every $\varepsilon>0$ and every
sufficiently large $N$ there is such a polynomial whose modulus lies
between $(1-\varepsilon)\sqrt N$ and $(1+\varepsilon)\sqrt N$ at every
point of the unit circle. Its proof gives no value for ``sufficiently
large''. We give one: for $0<\varepsilon<1$, every
$N\ge\exp(\exp(\exp(10^{12}\log^2(2/\varepsilon))))$ works. In particular,
for every $N\ge\exp(\exp(\exp(10^{12}\log^2 4)))$ the error can be taken to be
$2\exp(-10^{-6}\sqrt{\log\log\log N})$. The proof follows the original
construction and makes each step explicit. Two steps need new arguments:
a hypergraph matching lemma whose constants are explicit in the size of the
edges, replacing the Pippenger--Spencer theorem, and an exponential decay
estimate showing that the auxiliary polynomials need only logarithmically
many variables. The bound is far too large to be practical. Its point is
that the construction is effective, with every constant written down.
\end{abstract}

\noindent\textbf{Status and AI assistance.}
The proofs in this paper were written with the assistance of AI models and
were checked by separate AI proof audits. They have not been reviewed by a
human expert or formally verified. Novelty is qualified in
Section~\ref{sec:prior}.

\tableofcontents

\section{Introduction}

A real Littlewood polynomial of length $N$ has the form
$P(z)=\sum_{j=0}^{N-1}\epsilon_jz^j$ with every $\epsilon_j\in\{-1,1\}$.
By Parseval's identity the mean of $|P(z)|^2$ over the unit circle is $N$,
so $\sqrt N$ is the natural size of $|P|$. A family of such polynomials is
\emph{ultraflat} if $|P(z)|/\sqrt N$ tends to $1$ uniformly on the circle
as the length grows. OpenAI's preprint \cite{main} proved that real
Littlewood polynomials can be ultraflat at every sufficiently large length.
Its argument fixes many auxiliary objects by qualitative choices, so it gives
no bound on how large the length must be for a given accuracy.

This paper makes that argument explicit. Theorem~\ref{thm:main} says that
accuracy $\varepsilon$ is reached at every length
$N\ge\exp(\exp(\exp(10^{12}\log^2(2/\varepsilon))))$, and
Corollary~\ref{cor:rate} restates this as an error of
$2\exp(-10^{-6}\sqrt{\log\log\log N})$ at length $N$. Two steps need new
arguments. The original construction packs short arcs on the circle using
the Pippenger--Spencer theorem \cite{pippengerspencer} on matchings in
hypergraphs, in the qualitative form of Alon and Yuster \cite{alonyuster};
Section~\ref{sec:matching} proves a matching lemma whose constants are
explicit in the size of the edges. The original construction
also builds an auxiliary polynomial by an iteration whose length was chosen
qualitatively; Lemma~\ref{lem:saturation} shows that its error decays
exponentially, so logarithmically many steps suffice. The remaining steps
re-derive the original estimates with numerical constants.

The bound is far too large to be practical, and no polynomial at these
lengths has been computed. The three nested exponentials come from the
packing step, where the size of the required prime grows exponentially in
the number of arcs.

\subsection{Prior work and novelty}\label{sec:prior}

For complex coefficients of modulus one, Kahane \cite{kahane} constructed
ultraflat polynomials with an explicit rate, and Bombieri and Bourgain
\cite{bombieribourgain} improved it. For real signs, Balister, Bollob\'as,
Morris, Sahasrabudhe and Tiba \cite{bbmst} proved that the modulus can be
kept between two constant multiples of $\sqrt N$, and \cite{main} proved
ultraflatness with no rate. Erd\'elyi \cite{erdelyi} states a lower bound
for the maximum modulus that is compatible with the rate here. The
construction in this paper is that of \cite{main,predecessor}; the
contribution is the explicit bound and the two arguments above.

Quantitative versions of the matching theorem exist. Alon, Kim and Spencer
\cite{alonkimspencer} and Kang, K\"uhn, Methuku and Osthus \cite{kkmo}
bound the number of uncovered vertices with explicit exponents. As far as we
found, their constants are not given explicitly as functions of the edge
size, which is what is needed here because the edge size grows with the
accuracy. The lemma in Section~\ref{sec:matching} is a simple nibble argument
written to give such constants; we do not claim it improves on those results.

On October 10, 2026 we searched arXiv (queries on ``ultraflat'' and on
``Littlewood polynomials'', through submissions of October 8), the openai/math
preprint repository, and the web. We found no earlier explicit rate for real
Littlewood polynomials. This is a bounded search, not a proof that none
exists.

\section{Statement and dependencies}

Write $e(t)=\exp(2\pi i t)$ and $\T=\R/\Z$, with normalized Haar measure.
All logarithms are natural unless a base is displayed. Set
\[
 \Ethree(x)=\exp(\exp(\exp x)).
\]
A real Littlewood polynomial of length $N$ is
$P(z)=\sum_{j=0}^{N-1}\epsilon_jz^j$ with $\epsilon_j\in\{-1,1\}$.
Its normalized modulus error is
\[
 \mathcal D(P)=\sup_{|z|=1}
 \left|\frac{|P(z)|}{\sqrt N}-1\right|.
\]

\begin{theorem}\label{thm:main}
For every $0<\varepsilon<1$ and every integer
\[
 N\ge \Ethree\!\left(10^{12}\log^2\frac2\varepsilon\right),
\]
there is a real Littlewood polynomial $P_N$ of length $N$ with
$\mathcal D(P_N)\le\varepsilon$.
\end{theorem}

\begin{corollary}\label{cor:rate}
For every integer $N\ge\Ethree(10^{12}\log^2 4)$, there is a real
Littlewood polynomial of length $N$ such that
\[
 \mathcal D(P_N)\le
 2\exp\!\left(-10^{-6}\sqrt{\log\log\log N}\right).
\]
The signs may be chosen separately for each length.
\end{corollary}

\begin{proof}
Apply Theorem~\ref{thm:main} with
$\varepsilon=2\exp(-10^{-6}\sqrt{\log\log\log N})$.
The stated lower bound on $N$ makes $0<\varepsilon\le1/2$, and substitution
in the theorem's threshold gives exactly $N$.
\end{proof}

The starting architecture is the auxiliary-torus, signed-packing and phase
construction of \cite{main,predecessor}. We record its required quantitative
estimates below. The standard external discrepancy input is the partial
coloring theorem of Lovett--Meka \cite[Theorem 4]{lovettmeka}; Bertrand's
postulate supplies a prime in a prescribed factor-two interval. The matching
argument below is self-contained and replaces the qualitative
Pippenger--Spencer input in the source construction. No assertion about the
running time of a coefficient-generation algorithm is made.

\section{An explicit matching lemma}\label{sec:matching}

\begin{theorem}\label{thm:matching}
Let $r\ge2$, $0<p\le1/2$, $0<\eta\le1/4$, and let
\[
 q\ge \left(\frac{2^{20}r}{p\eta}\right)^{328r}.
\]
Suppose a finite simple $r$-uniform hypergraph has $n\le rq$ vertices and
a real reference degree $D\ge p^{r-1}q$. Put
$a_0=(1+p^{-(r-1)})q^{-1/8}$. If every vertex has degree in
$[(1-a_0)D,(1+a_0)D]$ and every pair-codegree is at most
$D/(p^{r-1}q)$, then a matching covers at least $(1-\eta)n$ vertices.
\end{theorem}

We first prove a concentration estimate with its numerical constant.

\begin{lemma}\label{lem:bernstein}
Suppose $X$ is a function of $m$ independent Bernoulli$(\xi)$ variables,
and changing one variable changes $X$ by at most $A$. With
$V=m\xi(1-\xi)A^2$, for $u>0$,
\[
 \Prob(|X-\E X|\ge u)
 \le2\exp\left(-\frac{u^2}{2(V+Au/3)}\right).
\]
The zero-variance case is interpreted deterministically.
\end{lemma}

\begin{proof}
Expose the variables successively. Each Doob difference is
$(Z-\xi)\Delta$, where coupling the same future variables gives
$|\Delta|\le A$. Its conditional variance is at most
$\xi(1-\xi)A^2$. If $Y$ is a centered difference with $|Y|\le A$, then
$j!\ge2\cdot3^{j-2}$ for $j\ge2$ gives
\[
 \E(e^{\lambda Y}\mid\text{past})
 \le\exp\left(
 \frac{\lambda^2\E(Y^2\mid\text{past})}{2(1-\lambda A/3)}\right),
 \qquad 0<\lambda<3/A.
\]
Iterate this inequality and take $\lambda=u/(V+Au/3)$. Apply the same
argument to $-X$ for the other tail.
\end{proof}

\begin{proof}[Proof of Theorem~\ref{thm:matching}]
Put $k=r-1$ and define
\begin{align*}
 L&=\log(8/\eta),&
 \theta&=\frac{\eta}{2048rL},&
 T_0&=\ceil{4L/\theta},\\
 S&=(\eta/8)^{10r},&
 E_0&=(8/\eta)^{40r},&
 s&=\frac{\theta}{256LE_0},\\
 B_0&=2^{32}r^3L^3\eta^{-1}p^{-2k}(8/\eta)^{100r}.
\end{align*}
It suffices initially to assume $q\ge Q_*$, where
\begin{equation}\label{eq:Qstar}
 Q_* =\max\left\{
 (32p^{-k}E_0)^8,\quad
 640r^2LE_0p^{-k}S^{-1},\quad
 p^{-k}S^{-1},\quad
 (32rT_0B_0)^2\right\}.
\end{equation}
We verify the simpler stated threshold at the end.

Let $c$ denote the maximum pair-codegree in the original hypergraph. The
case $n=0$ is immediate. Otherwise the hypotheses below ensure positive
degrees, so $c>0$. Define deterministic reference degrees by
\[
 D_0=D,\qquad \xi_t=\theta/D_t,\qquad
 \rho_t=(1-\xi_t)^{(r-1)D_t},\qquad D_{t+1}=D_t\rho_t.
\]
At round $t$, mark each current edge independently with probability $\xi_t$.
Accept each marked edge that meets no other marked edge, and delete every
vertex in any marked edge. Continue on the induced surviving hypergraph.
Accepted edges from all rounds form a matching.

The inequalities $\theta T_0\le5L$ and $\log(1-x)\ge-2x$ for
$0\le x\le1/2$ imply, inductively,
\begin{equation}\label{eq:reference-degree}
 D_t\ge D\exp(-2r\theta t)\ge DS\ge1,
 \qquad \rho_t\ge e^{-2r\theta}\ge1/2.
\end{equation}
There is no circularity: the bound at time $t$ gives $\xi_t\le\theta<1/2$
and hence the bound at time $t+1$. Set
\[
 \kappa=(p^kqS)^{-1}.
\]
Pair-codegrees do not increase on taking induced subhypergraphs, and
\eqref{eq:Qstar} gives
\begin{equation}\label{eq:initial-smallness}
 \frac c{D_t}\le\kappa\le\frac1{640r^2LE_0},
 \qquad a_0\le\frac1{16E_0}.
\end{equation}

\paragraph{Conditional degree expectation.}
Suppose all current degrees are in $(1\pm a_t)D_t$, with $a_t\le1/4$.
Condition on a vertex $v$ surviving this round. This forces exactly the
edges incident to $v$ to have mark zero; every other mark remains an
independent Bernoulli$(\xi_t)$ variable.
For $e\ni v$, let $m_e$ count edges not containing $v$ that meet
$e\setminus\{v\}$. Removing the $v$-incident edges and subtracting pairwise
overlaps gives
\[
 (r-1)(1-a_t)D_t-\frac{r(r-1)}2c
 \le m_e\le(r-1)(1+a_t)D_t.
\]
The conditional survival probability of $e$ is $(1-\xi_t)^{m_e}$.
Consequently the conditional mean $\mu_v$ of the new degree satisfies
\[
 (1-a_t)e^{-z_t}\le\frac{\mu_v}{D_{t+1}}
 \le(1+a_t)e^{z_t},\qquad
 z_t=2\theta(ra_t+r^2c/D_t)\le1.
\]
Using $e^z\le1+2z$ on $[0,1]$ gives
\begin{equation}\label{eq:mean-degree}
 \left|\frac{\mu_v}{D_{t+1}}-1\right|
 \le(1+8r\theta)a_t+8r^2\theta\kappa.
\end{equation}

\paragraph{Conditional concentration.}
Changing an unconditioned mark on an edge $f$ changes the new degree of
$v$ by at most $\sum_{u\in f}\codeg(v,u)\le rc$.
There are at most $2nD_t/r$ marks. Lemma~\ref{lem:bernstein} applies with
$A=rc$ and $V\le2n\theta rc^2$. Since $D_{t+1}\ge D_t/2$,
\begin{align}
 &\Prob\bigl(|d_{t+1}(v)-\mu_v|>sD_{t+1}
       \mid v\text{ survives},\,\text{current hypergraph}\bigr)\notag\\
 &\quad\le2\exp\left[-\min\left\{
  \frac{s^2D_t^2}{32n\theta rc^2},\frac{sD_t}{8rc}\right\}\right]
 \le2e^{-q/B_0}.                                      \label{eq:degree-tail}
\end{align}
Indeed $n\le rq$ and $D_t/c\ge p^kqS$, and
\[
 \frac{32\theta r^2}{s^2p^{2k}S^2}=B_0,
 \qquad \frac{8r}{sp^kS}\le B_0.
\]
Multiplying the conditional probability by the probability of survival
also bounds the probability of a surviving vertex with this deviation.

\paragraph{Closing the induction.}
Define
\[
 a_{t+1}=(1+8r\theta)a_t+8r^2\theta\kappa+s.
\]
For every $t\le T_0$, this deterministic recurrence gives
\[
 a_t\le E_0\left(a_0+40r^2L\kappa+\frac{5Ls}{\theta}\right)
 \le\frac1{16}+\frac1{16}+\frac5{256}<\frac14.
\]
Thus the hypotheses needed at each good round persist. A union bound over
at most $nT_0$ vertex-round positions shows that the probability of any
degree failure is at most $2rqT_0e^{-q/B_0}\le1/8$.
For the last inequality, \eqref{eq:Qstar} gives
$q/B_0\ge32rT_0\sqrt q$, whereas
\[
 \log(16rqT_0)\le(16rT_0+1)\sqrt q\le32rT_0\sqrt q.
\]
Here $\log q\le\sqrt q$ for $q\ge1$.

\paragraph{Survivors and collision losses.}
Stop the process on the first degree failure. On a good current history,
each vertex survives the next round with probability at most
$e^{-3\theta/4}$. Let $G_t$ be the event that all degree conditions hold
through time $t$, and let $V_t$ be the surviving vertex set. Retaining the
good-history indicators in conditional expectation gives
\[
 \E(|V_{t+1}|\mathbf1_{G_{t+1}})
 \le e^{-3\theta/4}\E(|V_t|\mathbf1_{G_t}).
\]
It follows that
\[
 \E(|V_{T_0}|\mathbf1_{G_{T_0}})\le n(\eta/8)^3,
 \qquad
 \Prob(G_{T_0},\ |V_{T_0}|>\eta n/2)\le\eta^2/256\le1/4096.
\]
Every removed vertex that is not covered by an accepted edge lies in a
marked edge having another marked neighbor. The number of such vertices
is at most $r$ times the number of ordered intersecting pairs of distinct
marked edges. There are at most $2nD_t/r$ edges and at most $2rD_t$
neighbors per edge, so expected waste in one executed round is at most
$4nr\theta^2$. Expected total waste is therefore at most
$20nr\theta L=20\eta n/2048$. The probability that it exceeds $\eta n/2$
is less than $1/32$. Together the three failure bounds sum to less than
one. Some outcome has at most $\eta n$ uncovered vertices.

\paragraph{The stated threshold.}
Set $W=2^{20}r/(p\eta)\ge2^{24}$. Each of $r,p^{-1},\eta^{-1},L$ is at
most $W$, and $8/\eta\le W$, $2^{32}\le W^2$. Hence
\[
 T_0\le W^5,\quad E_0\le W^{40r},\quad S^{-1}\le W^{10r},
 \quad B_0\le W^{102r+7}.
\]
The four terms in \eqref{eq:Qstar} are at most
$W^{328r}$, $W^{51r+3}$, $W^{11r-1}$ and $W^{204r+28}$, respectively.
For $r\ge2$ all are at most $W^{328r}$, completing the proof.
\end{proof}

\section{Quantitative auxiliary data}\label{sec:auxiliary}

\subsection{The discrepancy input}

We use the following explicit form of \cite[Lemma 3.2]{main}.

\begin{lemma}\label{lem:rounding}
Let $Z_0,\ldots,Z_{n-1}\in\C$ satisfy $|Z_j|\le1$. Write
$Z_j=d_j|Z_j|$ with $|d_j|=1$, taking $d_j=1$ when $Z_j=0$, and put
$\mu=\frac12\sum_j(1-|Z_j|)$. There are $\eta_j\in\{d_j,-d_j\}$ such that
\begin{equation}\label{eq:rounding}
 \sup_{t\in\T}\left|\sum_{j=0}^{n-1}(\eta_j-Z_j)e(jt)\right|
 \le C_R\left(1+\sqrt{\mu\log(80n/\mu)}\right),\qquad C_R=1024.
\end{equation}
If $\mu=0$ the error can be zero. For real inputs the outputs are real
signs. The same assertion holds on any consecutive interval of frequencies.
\end{lemma}

\begin{proof}[Derivation of the numerical constant]
The Lovett--Meka partial coloring theorem \cite[Theorem 4]{lovettmeka}
allows arbitrary vectors and arbitrary initial points in the cube. Taking
the tolerance to zero by compactness gives an exact partial coloring that
fixes at least half the active coordinates. For a real $R\times s$ matrix
with entries in $[-1,1]$ and $R\ge s$, use the threshold
$4\sqrt{\log(16R/u)}$ when $u$ coordinates remain active. The total
discrepancy in a complete coloring is bounded by
\[
 4\sqrt{sL_0}\sum_{h\ge0}2^{-h/2}\sqrt{1+h/4},
 \qquad L_0=\log(16R/s).
\]
The series is less than $7$, using $2^{-1/2}<3/4$ and
$\sqrt{1+h/4}\le1+h/4$. Since $L_0\le4\log(2R/s)$, a valid complete
matrix discrepancy constant is $64$.

For completeness, set $p_j=(1-|Z_j|)/2$ and sample real and imaginary
parts of $d_je(jt)$ on the $20n$ equally spaced points of $\T$. This gives
$R=40n$ real rows. First round $p_j$ down to a dyadic mesh of size $2^{-J}$
with $n2^{-J}\le1$, costing at most $1$ in row norm. At stage $h$, color
the coordinates with odd numerator on that mesh, reverse all signs if
necessary to keep their mass nonincreasing, and move them by $2^{-h}$.
Their number $s_h$ is at most $\min(n,2^h\mu)$. Monotonicity of
$s\log(80n/s)$ and the geometric sum
$\sum_{h\ge1}2^{-h/2}=1+\sqrt2<5/2$ give row error at most
$1+160\sqrt{\mu\log(80n/\mu)}$.
Passing to complex values costs $2\sqrt2$. The polynomial derivative
bound used in \cite[Lemma 3.2]{main} costs at most
$(1-\pi\exp(1)/20)^{-1}<2$ in passing from the grid to the circle.
These factors are bounded by the constant in \eqref{eq:rounding}.
For real inputs all $d_j$ are real signs. A unimodular monomial handles
shifted frequency intervals.
\end{proof}

\subsection{Exponential saturation of the auxiliary recursion}

Let
\[
 p_0=0,\qquad
 p_j(y_1,\ldots,y_j)=p_{j-1}
       +\frac12(1-p_{j-1}^2)\cos(2\pi y_j).
\]
These are real trigonometric polynomials, $|p_j|\le1$, and their means
are zero.

\begin{lemma}\label{lem:saturation}
For every integer $d\ge0$,
\[
 1-\int_{\T^d}p_d^2\le(7/8)^{d/2}.
\]
\end{lemma}

\begin{proof}
Let $H_0(x)=\sqrt{1-x^2}$. Fix $|x|<1$, write $u=1-x^2$, and pair
$c=\cos(2\pi y)$ with $-c$. Put
$A=1-uc^2/4$ and $D_0=1-c^2/2+uc^2/4$. The square of the paired mean
of $H_0(x+uc/2)/\sqrt u$ is
\[
 \frac{A+\sqrt{A^2-x^2c^2}}2.
\]
Here $D_0\ge1/2$ and
$D_0^2-(A^2-x^2c^2)=(1-u)c^4/4\ge0$. The displayed expression is
therefore at most $(A+D_0)/2=1-c^2/4$. Jensen's inequality and
$\E c^2=1/2$ give
\[
 \E_y H_0\left(x+\tfrac12(1-x^2)\cos(2\pi y)\right)
 \le\sqrt{7/8}\,H_0(x).
\]
The endpoints $x=\pm1$ are immediate. Iterate from $p_0=0$ and use
$H_0^2\le H_0$.
\end{proof}

\subsection{A closed auxiliary envelope}

Fix $0<\delta<1/100$. The following symbols concern only this section;
in particular $T$ is the spreading parameter, not the number of matching
rounds. Define
\begin{align}
 d&=\max\left(2,\ceil{\frac{2\log(1/\delta)}{\log(8/7)}}\right),
 &H_0&=2^{d-1},& k_0&=2^{d+1}-2,\notag\\
 K&=3^d2^{d(d-1)/2},& B_1&=2H_0+2,&G_0&=dH_0B_1^{d-1},\notag\\
 P_0&=d^2H_0^2K,&\tau&=\frac{\delta}{4KH_0},
 &\alpha&=\frac{\delta}{4K\sqrt K},                  \label{eq:aux-data}\\
 x_0&=\frac{\alpha}{64C_R},&q_0&=x_0^4,&h_0&=q_0/12,\notag\\
 K_1&=64,&K_3&=2^{26},&M_2&=1+2K_3h_0^{-3},\notag\\
 E_1&=(2K_1+6)h_0^{-2}+3h_0^{-3},&&
 \beta_*&=\frac{\tau 2^{-2k_0}}{G_0P_0^{K-1}}.\notag
\end{align}
Finally put
\begin{align}
 R_0&=\max\left\{h_0^{-1},\frac{1+2\pi M_2}{q_0},\frac{E_1}{q_0},
                 1+(4C_R/\alpha)^2,\frac{4K}{\delta}\right\},
 &T&=\ceil{R_0/\beta_*}.                             \label{eq:spreading-T}
\end{align}

\begin{proposition}\label{prop:auxiliary}
There is a real trigonometric polynomial
$F(y)=\sum_{a\in A}c_ae(a\cdot y)$ on $\T^m$ and a vector $v\in\R^m$
such that $A$ consists of at least two nonparallel sign pairs,
$\lambda_a=a\cdot v\ne0$, $w_a=|\lambda_a|$, and
\begin{align}
 \norm F_\infty&\le1+\delta,\notag\\
 \frac{(1-\delta)^3}{(1+\delta)^3}
 &\le\sum_{a\in A}w_a\le\frac1{1+\delta},             \label{eq:aux-conclusions}\\
 1+\delta&\le\frac{|c_a|^2}{w_a}
      \le\frac{(1+\delta)^3}{(1-\delta)^2},\qquad
 1\le\frac{|c_a|}{\sqrt{w_a}}\le1+4\delta<2.\notag
\end{align}
The parameters obey
\begin{align}
 m&\le d+K/2,& |A|&\le K\bigl(T(P_0+1)+1\bigr)^K,\notag\\
 \max_{a\in A}\norm a_\infty
 &\le\max\{H_0,T(P_0+1)/2\},&
 \min_{a\in A}w_a&\ge\frac1{\tau T^K(1+\delta)^3}.    \label{eq:aux-envelope}
\end{align}
All these data are fixed before the final length $N$ is chosen.
\end{proposition}

\begin{proof}
We quantify the construction in \cite[Section 4]{main}.
Write $p_d=\sum_{s\in\mathcal S}b_se(s\cdot y)$.
Every coefficient is an integer multiple of $2^{-k_0}$, since the
denominator exponents satisfy $k_j=2k_{j-1}+2$. Coordinate degrees satisfy
$|s_i|\le2^{d-i}\le H_0$, and consequently $|\mathcal S|\le K$.
Every frequency's last nonzero coordinate is $\pm1$, by induction in
the displayed recursion. Thus distinct sign pairs are nonparallel.
Lemma~\ref{lem:saturation} and Parseval give
\[
 1-\delta\le\sum_s|b_s|^2\le1,\qquad
 \sum_s|b_s|\le\sqrt K,\qquad |b_s|\ge2^{-k_0}.
\]

Set $\gamma=(1,B_1,\ldots,B_1^{d-1})$ and
$\rho_s=|s\cdot\gamma|$. Leading-digit domination gives
$1\le\rho_s\le G_0$. Choose representatives $s_1,\ldots,s_D$, where
$2\le D\le K/2$. For each $j$ choose an integer vector
$W_j^0\in s_j^\perp$ that is not orthogonal to another representative.
Here is a bounded choice. If the last nonzero coordinate of $s_j$ is
at $i_j$, use the integer basis
$e_i-(s_j)_i(s_j)_{i_j}^{-1}e_{i_j}$, $i\ne i_j$, and coefficients
in $\{0,\ldots,D-1\}$. Each of the $D-1$ forbidden hyperplanes contains
at most $D^{d-2}$ grid points, fewer in total than $D^{d-1}$.
This gives
\[
 \norm{W_j^0}_\infty\le dH_0K,
 \qquad 1\le|s_\ell\cdot W_j^0|\le P_0\quad(\ell\ne j).
\]
Put
\[
 Q_j=\prod_{\ell\ne j}|s_j\cdot W_\ell^0|,
 \quad u_j=\frac{|b_{s_j}|^2}{\rho_{s_j}Q_j},
 \quad W_j=W_j^0+\tau\frac{u_j}{(s_j)_{i_j}}e_{i_j}.
\]
For $\beta_{s,j}=s\cdot W_j$, each off-diagonal factor has relative
perturbation at most $\tau H_0\le\delta/(4K)$, and the diagonal factor
is exactly $\tau u_j$ in absolute value. Hence
\begin{equation}\label{eq:beta-product}
 (1-\delta)\tau\frac{|b_s|^2}{\rho_s}
 \le\prod_{j=1}^D|\beta_{s,j}|
 \le(1+\delta)\tau\frac{|b_s|^2}{\rho_s},
 \qquad \beta_*\le|\beta_{s,j}|\le P_0+1.
\end{equation}

We next record why \eqref{eq:spreading-T} suffices for every spreading
factor. Take $\kappa_0(x)=C\exp(-1/(1-x^2))$ for $|x|<1$, zero otherwise,
normalized to have integral one. Its constant satisfies $C<4$, and direct
differentiation gives
$\norm{\kappa_0'}_1\le K_1$ and $\norm{\kappa_0'''}_1\le K_3$.
For example the third derivative is bounded by
$4\cdot140(1-x^2)^{-6}\exp(-1/(1-x^2))$; the bound
$2\cdot4\cdot140\cdot6^6<2^{26}$ suffices.
Let $g=\mathbf1_{[-1/2+2h_0,\,1/2-2h_0]}*\kappa_{0,h_0}$.
Then $0\le g\le1$, its support stays $h_0$ from the endpoints,
\[
 \int g^2\ge1-q_0/2,\quad
 \norm{g'}_1\le2,\quad\norm{g''}_1\le2K_1/h_0,
 \quad\norm{g^{(4)}}_1\le2K_3/h_0^3,
 \quad\int\xi^2|\widehat g(\xi)|\,d\xi\le M_2.
\]
For a nonzero $\beta$ in the finite list, set $a=T|\beta|$,
$L_\beta=2\floor{a/2}+1$, and
$I_\beta(u)=\int g(y)e(T\beta y^2/2-uy)\,dy$.
The exact quadratic Fourier identity gives
\[
 \sqrt{L_\beta}|I_\beta(u)|
 \le\sqrt{1+1/a}(1+\pi M_2/a)
 \le1+(1+2\pi M_2)/a\le1+q_0.
\]
For $|u|>\floor{a/2}$, put $v_u=|u|-a/2>0$.
Twice integrating by parts yields
\[
 |I_\beta(u)|\le\frac1{4\pi^2}
 \left(\frac{\norm{g''}_1}{(v_u+ah_0)^2}
 +\frac{3a\norm{g'}_1}{(v_u+ah_0)^3}
 +\frac{3a^2\norm g_1}{(v_u+ah_0)^4}\right).
\]
Since $ah_0\ge1$, the two-sided sum of $(v_u+ah_0)^{-j}$ is at most
$4(ah_0)^{1-j}$ for $j\ge2$. The exterior coefficient sum is therefore
at most $E_1/a\le q_0$. Parseval gives retained energy at least $1-2q_0$.
Apply Lemma~\ref{lem:rounding} to
$Z_u=\sqrt{L_\beta}I_\beta(u)/(1+q_0)$. Its normalized defect is at
most $2q_0$, so the resulting equal-magnitude polynomial $f_\beta$ obeys
\[
 \norm{f_\beta-g(\cdot)e(T\beta(\cdot)^2/2)}_\infty
 \le C_RL_\beta^{-1/2}
      +C_R\sqrt{2q_0\log(40/q_0)}+2q_0<\alpha.
\]
Indeed the three terms are bounded respectively by
$\alpha/4$, $\alpha/128$, and $\alpha/4$ from
\eqref{eq:aux-data}--\eqref{eq:spreading-T}. For the middle term one may
use $x_0\le1/64$, $\log80<5$, $\log(1/x_0)\le1/x_0$, and
$x_0\sqrt{2(\log80+4\log(1/x_0))}<1/2$.
Choose $f_{-\beta}=\overline{f_\beta}$.

Define, on $\T^{d+D}$,
\[
 F(y,z)=\sum_{s\in\mathcal S}b_se(s\cdot y)
                       \prod_{j=1}^D f_{\beta_{s,j}}(z_j).
\]
Replacing every factor by its target gives
\[
 \left(\prod_jg(z_j)\right)
 p_d\left(y+\frac T2\sum_jW_jz_j^2\right),
\]
of norm at most one. Telescoping the products costs at most
$\sqrt K D\alpha(1+\alpha)^{D-1}<\delta$.
Each box over $s$ has $n_s=\prod_jL_{\beta_{s,j}}$ frequencies, all with
coefficient modulus $|b_s|/\sqrt{n_s}$. Moreover
\[
 (1-\delta)T^D\prod_j|\beta_{s,j}|
 \le n_s\le(1+\delta)T^D\prod_j|\beta_{s,j}|,
\]
since every factor $L_\beta/(T|\beta|)$ differs from one by at most
$\delta/(4K)$. Set
\[
 v=\left(\frac{\gamma}{\tau T^D(1+\delta)^3},0\right).
\]
The weight on the box over $s$ is
$\rho_s/(\tau T^D(1+\delta)^3)$. Combining the bounds for $n_s$ with
\eqref{eq:beta-product} proves \eqref{eq:aux-conclusions} and
\eqref{eq:aux-envelope}. Nonparallelism follows by projecting a putative
parallel pair to its original $s$ coordinates.
\end{proof}

\begin{remark}
There is no necessary unbounded search for coefficient phase precision.
Let $L_{\max}=T(P_0+1)+1$ and
$R_{\rm ph}=\ceil{64\sqrt{L_{\max}}/\alpha}$.
Rounding each positive-$\beta$ unit coefficient to a nearest
$R_{\rm ph}$-th root of unity changes the normalized factor by at most
$\alpha/16$. The preceding strict estimate has this much room.
Conjugating the negative-$\beta$ factors preserves all conclusions.
The later analytic bounds depend only on moduli and phase derivatives,
not on denominators of these arguments.
\end{remark}

\section{Quantitative signed packing}\label{sec:packing}

Let $a_1,\ldots,a_r\in\Z^m\setminus\{0\}$ be pairwise nonparallel,
$r\ge2$, and let $w_i>0$ satisfy $2\sum_iw_i<1$. We apply the following
construction to representatives of the sign pairs in
Proposition~\ref{prop:auxiliary}. Write
\begin{align*}
 \Delta&=1-2\sum_iw_i,&
 \alpha_i&=2w_i+\Delta/r,&
 M_0&=\ceil{4r(r+1)/\Delta}+1.
\end{align*}
Choose $L_i$ to be a nearest positive odd integer to $M_0\alpha_i$, and set
\[
 L=\sum_iL_i,\qquad p=\frac1{2L},\qquad
 k=r-1,\qquad \eta=\frac{\Delta}{16r}.
\]
Since $|L_i-M_0\alpha_i|\le1$ and $M_0\Delta/r\ge4(r+1)$,
\begin{equation}\label{eq:slot-slack}
 L_i-2w_iL\ge\frac{3M_0\Delta}{4r},\qquad
 1-\max_i\frac{2Lw_i}{L_i}\ge\frac{\Delta}{2r}=8\eta.
\end{equation}
For the auxiliary data, $\Delta>\delta/2$ and thus
$L\le64r^2/\delta$, $\eta\ge\delta/(32r)$, and $\eta\le1/32$.

Let $b\ge2$ be the rational rank of the row matrix. Let $A_0\ge1$ bound
the absolute values of one nonzero $b$-minor and one nonzero $2$-minor for
every pair of rows. Put
\[
 C_{
m mom}=(24kL+18k^2)p^{-2k}.
\]
Choose a prime $q$ strictly greater than all the following quantities:
\begin{equation}\label{eq:packing-threshold}
 \begin{gathered}
 4,\quad A_0,\quad
 [2r(C_{\rm mom}+3k(k-1)+1)]^8,\quad
 [(1+8L)/\eta]^8,\quad
 [2^{20}r/(p\eta)]^{2^{20}r^2}.
 \end{gathered}
\end{equation}
The last term is a convenient larger threshold than the one in
Theorem~\ref{thm:matching}.

\begin{proposition}\label{prop:packing}
For this $q$, there are an integer $1\le H\le q$ and points
$\theta_1,\ldots,\theta_H\in q^{-1}\Z^m/\Z^m$ such that the signed
closed arcs centered at $\pm a_i\cdot\theta_h$, of lengths $w_i/H$,
are pairwise disjoint. The arcs have pairwise separation at least
$6\eta/q$, and their distances from $0$ and $1/2$ are at least $\eta/q$.
Moreover
\[
 H\ge(1-\eta)^2\frac q{2L}\ge(1-2\eta)\frac q{2L}.
\]
\end{proposition}

\begin{proof}
We quantify the finite-field construction in
\cite[Section 3]{predecessor}. The choice $q>A_0$ preserves the full rank
$b$ and rank two for every row pair over $\mathbb F_q$.
Put $R=\floor{\sqrt q}$, $n_+=(q-1)/2$, and divide the positive sites
$\{1,\ldots,n_+\}$ into chunks of $R$ consecutive sites, with a possibly
shorter final chunk. Their number $K_c$ is at most $3\sqrt q/2$.
In each chunk independently shift a periodic pattern of consecutive slots
of lengths $L_1,\ldots,L_r$ by a uniform residue modulo $L$ and retain
only full slots. Each retained slot is independently oriented with a fair
sign. Since $L_i$ is odd its center is an integer site.
The number of type-$i$ slots in a chunk of $u$ sites differs from $u/L$
by at most $2$. Thus every realization has between
$n_+/L-2K_c$ and $n_+/L+2K_c$ slots of each type.

For $s\in\mathbb F_q$, let $|s|_q$ be the absolute value of its
representative in $[-n_+,n_+]$, and let $C_i(s)$ indicate that $s$ is an
oriented center of type $i$. A site is interior if it is at distance at
least $L$ from both ends of its chunk. For interior folded sites,
$\E C_i(s)=p=1/(2L)$. At most $2LK_c$ positive sites are noninterior.
An edge is a compatible tuple of one slot of each type, meaning that
$a_i\cdot y=s_i$ for all $i$ has a solution in $\mathbb F_q^m$.
Each tuple is included once, making a simple $r$-uniform hypergraph.

Fix a candidate center $s\ne0$ of type $i$ and condition on its presence.
This affects only the chunk of $|s|_q$. Exclude centers $s$ for which
some $j<\ell$, both different from $i$, satisfy
$a_j\pm a_\ell=c a_i$ with $c\ne0$ and $|cs|_q\le R$.
There are at most
\[
 2\binom{k}{2}(2R+1)\le3k(k-1)\sqrt q
\]
such centers for each type. Pairwise rank two excludes a zero constant
combination. For a nonexceptional $s$, sample $y$ uniformly from
$a_i\cdot y=s$. Each other evaluation is uniform; the probability that
one evaluation is zero, noninterior or in the conditioned chunk is at most
\[
 \frac{1+4LK_c+2R}{q}\le\frac{6L+3}{\sqrt q}.
\]
Two evaluations in the same folded chunk have sum or difference of
folded size at most $R$. A nonconstant combination is uniform, and each
sign costs at most $(2R+1)/q\le3/\sqrt q$. Constant within-copy
combinations were excluded above. Across two independent fiber samples,
the relevant evaluations are independent and uniform, including when the
row indices agree. For $u=1,2$ copies, a union bound therefore gives
\begin{equation}\label{eq:fiber-bad}
 F_u(q)\le\frac{uk(6L+3uk)}{\sqrt q}
\end{equation}
for failure to occupy distinct interior unconditioned chunks.

Define
\[
 Z_{i,s}=q^{-(m-1)}
       \sum_{a_i\cdot y=s}\prod_{j\ne i}C_j(a_j\cdot y).
\]
On good sample tuples the center events involve distinct chunks, so their
conditional probability is exactly $p^{uk}$. It follows that
\[
 |\E Z_{i,s}-p^k|\le F_1,
 \quad |\E Z_{i,s}^2-p^{2k}|\le F_2,
 \quad \E(Z_{i,s}-p^k)^2
      \le\frac{24kL+18k^2}{\sqrt q}.
\]
All expectations are conditional on the candidate vertex being present.
Chebyshev with relative error $q^{-1/8}$ bounds its failure probability
by $C_{\rm mom}q^{-1/4}$. There are at most $r(q-1)$ candidates. Together
with the exceptional centers, the expected number of deletions is at most
\[
 rC_{\rm mom}q^{3/4}+3rk(k-1)\sqrt q<q^{7/8},
\]
by \eqref{eq:packing-threshold}. Fix a realization with fewer than
$q^{7/8}$ deletions.

The evaluation map has $q^{m-b}$ preimages per compatible tuple. A
present vertex therefore has degree $q^{b-1}Z_{i,s}$, and fixing two
different types gives pair-codegree at most $q^{b-2}$. Same-type
pair-codegrees are zero. Deletion loses fewer than
$q^{7/8}q^{b-2}$ incident edges at each remaining vertex. With
$D=p^kq^{b-1}$, the remaining degrees consequently lie in
\[
 \bigl(1\pm(1+p^{-k})q^{-1/8}\bigr)D,
\]
and the maximum pair-codegree is at most $D/(p^kq)$.
There are at most $rq$ vertices and $D\ge p^kq$. All hypotheses of
Theorem~\ref{thm:matching} hold.

The remaining vertex count $v$ obeys
\[
 \frac vr\ge\frac q{2L}-\frac1{2L}-2K_c-q^{7/8}
       >(1-\eta)\frac q{2L}.
\]
Indeed its relative deficit is at most $(1+8L)q^{-1/8}<\eta$.
The matching supplied by Theorem~\ref{thm:matching} has the asserted size
$H$, and $H\le q$. Choose one generating vector $y_h$ for each edge and
put $\theta_h=y_h/q$. By \eqref{eq:slot-slack},
\[
 \frac{L_i}{q}-\frac{w_i}{H}
 \ge\frac{6\eta}{1-2\eta}\frac{L_i}{q}\ge\frac{6\eta}{q}.
\]
Thus each desired arc is inset by at least $3\eta/q$ from its slot
boundary. Slot interiors are disjoint, and the positive slots lie in
$[1/(2q),1/2]$. Reflection gives the signed statement and the claimed
boundary distances.
\end{proof}

\section{Uniform analytic estimates}\label{sec:analytic}

We give a numerical stationary-phase certificate before using it in the
wave construction. These estimates quantify \cite[Sections 2 and 5]{main}.

\begin{lemma}\label{lem:stationary}
Let $I$ have length $\ell\le1$, and let $\phi$ be real $C^7$ on a
neighborhood of $I$, with fixed sign of curvature. Suppose
\[
 0<b\le\min(1,\inf_I|\phi''|),\qquad
 B\ge\max\left(1,\max_{2\le j\le7}\norm{\phi^{(j)}}_\infty\right).
\]
Let $f\in C^4$ have support at distance at least $h>0$ from the endpoints,
and let $F_j\ge\norm{f^{(j)}}_\infty$ for $0\le j\le4$. Define
\begin{align*}
 a_0&=B,&
 a_k&=\frac{B+\sum_{j=1}^{k-1}\binom{k}{j}a_ja_{k-j}}{2b}
       &&(1\le k\le5),\\
 Y_k&=\ell a_k+ka_{k-1}&&(1\le k\le5),\\
 R_0&=B/b,&
 R_k&=(B/b)\sum_{j=0}^{k-1}\binom{k}{j}R_jY_{k-j+1}
       &&(1\le k\le4),\\
 V_{0,k}&=\sum_{j=0}^k\binom{k}{j}F_jR_{k-j}&&(0\le k\le4),\\
 V_{a,k}&=\sum_{j=0}^k\binom{k}{j}R_jV_{a-1,k-j+1}
       &&(1\le a\le4, 0\le k\le4-a),\\
 C_M&=4\ell\bigl(F_0+(B/b)V_{4,0}\bigr),&
 C_{NS}&=\ell\left(\frac{F_1}{bh}+
                          \frac{F_0B}{b^2h^2}\right).
\end{align*}
For every real $x$ and $N\ge1$, the scaled integral
$\sqrt N\int_I f(t)e(N(\phi(t)-xt))\,dt$ differs from its usual
stationary-phase main term by at most $C_M/N+C_{NS}/\sqrt N$.
If $\phi'(u)=x$ for $u\in I$, that term is
\[
 \frac{f(u)}{\sqrt{|\phi''(u)|}}
 e\bigl(N(\phi(u)-xu)+\sgn(\phi'')/8\bigr);
\]
otherwise it is zero.
\end{lemma}

\begin{proof}
For $x=\phi'(u)$ put $d=\sgn(\phi'')$ and
\[
 A(t,u)=2d\int_0^1(1-s)\phi''(u+s(t-u))\,ds,
 \quad a(t,u)=\sqrt{A(t,u)},\quad y=(t-u)a(t,u).
\]
Then $b\le A\le B$, its first five $t$-derivatives are bounded by $B$,
and the phase becomes $\phi(u)-xu+dy^2/2$ exactly. Differentiating
$a^2=A$ and then $y=(t-u)a$ gives the bounds $a_k,Y_k$ above.
Also $b/\sqrt B\le y'\le B/\sqrt b$, so both $y'$ and $1/y'$ are
bounded by $B/b$. Writing $r=1/y'$, differentiation of $ry'=1$ gives
the $R_k$ bounds. The transformed amplitude $g(y)=f(t(y))r(t(y))$
satisfies $\norm g_1\le\ell F_0$ and
$\norm{g^{(4)}}_1\le\ell(B/b)V_{4,0}$, since $d/dy=r\,d/dt$.
It extends by zero as a compactly supported $C^4$ function.

The exact quadratic Fourier identity, obtained for instance with a
Gaussian regulator, gives error at most
$\pi N^{-1}\int\xi^2|\widehat g(\xi)|\,d\xi$.
Four integrations by parts show
\[
 \int\xi^2|\widehat g(\xi)|\,d\xi
 \le\tfrac23\norm g_1+\frac2{(2\pi)^4}\norm{g^{(4)}}_1
 \le\norm g_1+\norm{g^{(4)}}_1.
\]
This gives $C_M/N$. For $x$ outside $\phi'(I)$, the derivative
$\phi'-x$ has modulus at least $bh$ on the support. A single integration
by parts gives $C_{NS}/\sqrt N$. Endpoints of the derivative range are
covered by the first argument.
\end{proof}

\subsection{Main waves and their coherence}

Apply Proposition~\ref{prop:packing} to the sign-pair representatives of
$A$ from Proposition~\ref{prop:auxiliary}. The centers can be written
$t^*_{a,h}=-a\cdot\theta_h\pmod1$, since both signs are packed.
Write $m_a=|c_a|/\sqrt{w_a}\in[1,1+4\delta]$ and set
\begin{equation}\label{eq:cutoff-parameters}
 J_h=\left[\frac{h-1}H+\frac{\delta}{16H},
             \frac hH-\frac{\delta}{16H}\right],\quad
 e_0=\frac{\delta}{32H},\quad
 \sigma=\frac{\delta}{24H|A|}.
\end{equation}
Use the monotone $C^6$ cutoff
\[
 \zeta(s)=12012\int_0^s t^6(1-t)^6\,dt\quad(0\le s\le1),
\]
extended by zero and one. For distance $d_x$ to the nearest endpoint of
$J_h$, set $\chi_h=\sigma$ for $d_x\le2e_0$, interpolate to one using
$\zeta((d_x-2e_0)/e_0)$ on $[2e_0,3e_0]$, and keep it one thereafter.
Set $G_h=0$ for $d_x\le e_0/2$, interpolate to one on $[e_0/2,e_0]$
using $\zeta$, and keep it one thereafter. Both are constant near the
midpoint, so this definition introduces no singularity there. They satisfy
\[
 \sigma\le\chi_h\le1,\qquad
 \int_{J_h}\chi_h^2\ge\frac{1-5\delta/16}{H}
                         >\frac{1-\delta}{H}.
\]

Let $Q_h''=\chi_h^2$ with $Q_h=Q_h'=0$ at the midpoint. Define
\begin{align}
 T_{a,h}(x)&=t^*_{a,h}-\lambda_a Q_h'(x),&
 X_{a,h}&=T_{a,h}^{-1},\notag\\
 \psi_{a,h}(t)&=(t-t^*_{a,h})X_{a,h}(t)
                   +\lambda_aQ_h(X_{a,h}(t)).        \label{eq:main-phase}
\end{align}
The image $I_{a,h}=T_{a,h}(J_h)$ is strictly inside its packed arc.
The identities
\[
 \psi_{a,h}'=X_{a,h},\qquad
 \psi_{a,h}''=-\frac1{\lambda_a\chi_h(X_{a,h})^2}
\]
follow by differentiation. On $I_{a,h}$ define
\begin{equation}\label{eq:main-BN}
 B_N(t)=\frac{c_a}{\sqrt{w_a}}
             e\bigl(\sgn(\lambda_a)/8+N\psi_{a,h}(t)\bigr).
\end{equation}
The opposite frequency gives the conjugate reflection. Main intervals
have total length at least $(1-\delta)\sum_aw_a$.

Apply Lemma~\ref{lem:stationary} with the cutoff $G_h(X_{a,h})$ and
$x=k/N$. Its leading scaled contribution is, for $x\in J_h$,
\[
 c_aG_h(x)\chi_h(x)
       e\bigl(-kt^*_{a,h}+N\lambda_aQ_h(x)\bigr).
\]
At most one $J_h$ contains $x$. Summing over $a$ for that $h$ gives
\begin{equation}\label{eq:coherence}
 G_h(x)\chi_h(x)F\bigl(k\theta_h+NvQ_h(x)\bigr),
\end{equation}
whose modulus is at most $1+\delta$. The congruence defining $t^*$ is
legitimate here because $k$ is an integer.

The remaining two endpoint pieces lie where $\chi_h=\sigma$ and have
constant curvature of modulus $1/(w_a\sigma^2)$. The primitive of
$e(\pm y^2/2)$ has differences bounded by $3$ on arbitrary intervals;
split at $\pm1$ and integrate the tails by parts to see this. Including
amplitude variation, each endpoint contributes at most $6|c_a|\sigma$
to a scaled Fourier coefficient. The sum over both endpoints and all
$(a,h)$ is at most $12H\sigma\sum_a|c_a|\le\delta$.

\subsection{Joining the gaps with numerical separation}

On $[0,1/2]$, let the gaps have lengths $\ell_j$, number $M=rH+1$, and
total length $L_g$. Proposition~\ref{prop:packing} and the total main
length give
\begin{equation}\label{eq:gap-lengths}
 \ell_j\ge\eta/q,\qquad M\le rq+1,\qquad 0<L_g\le3.5\delta.
\end{equation}
Prescribe endpoint derivatives $x_z$ from the adjacent main waves, and
use $x_0=x_{1/2}=1/8$ at the boundary. Assign to gap $j$ an interval
$[P_j,R_j]$ of length $\ell_j/(4L_g)$. Arrange these successively starting
at $5/16$, separated by spaces $s_g=1/[8(M+1)]$.
They have total length $1/4$, total spacing less than $1/8$, and lie
below $11/16$.

There are at most $4M^2$ forbidden common translations that put one of
the $2M$ assigned endpoints on one of at most $2M$ prescribed values.
Take the $4M^2+1$ translations
\[
 \frac{j}{256M^2},\qquad 0\le j\le4M^2.
\]
Each forbidden value is within $1/(768M^2)$ of at most one candidate.
One candidate remains, and after that translation
\begin{equation}\label{eq:avoidance}
 |P_j-x_z|, |R_j-x_z|\ge\frac1{768M^2}
 \quad\text{for every prescribed }x_z.
\end{equation}
The intervals still lie in $(1/4,3/4)$. Their closed $\rho$-neighborhoods
are disjoint for $\rho=1/[32(M+1)]$.

On each gap choose a continuous piecewise linear leading derivative
interpolating
$x_{\rm left}\longrightarrow P_j\longrightarrow R_j
\longrightarrow x_{\rm right}$. Give its outer pieces lengths
\begin{equation}\label{eq:outer-length}
 u_j=\min\left\{\frac{\ell_j}{8},
             |P_j-x_{\rm left}|\left(\frac{\delta}{400M}\right)^2\right\}
\end{equation}
and the analogous expression on the right. Their nonzero curvatures
$\Lambda$ satisfy
$\sum_{\rm outer}|\Lambda|^{-1/2}\le\delta/200$ on the positive half.
The middle curvature is positive and at least $1/(4L_g)$.
Integrate to obtain a $C^1$, piecewise quadratic phase $\phi_j$.

Interpolate the endpoint moduli linearly by $m_j$, prescribing modulus
one at $0$ and $1/2$. Choose an affine $\gamma_{j,N}$ with both endpoint
values in $[0,1)$ so that
\[
 B_N(t)=m_j(t)e\bigl(N\phi_j(t)+\gamma_{j,N}(t)\bigr)
\]
matches the two required endpoint values. This is always possible by
choosing representatives of the two phase congruences. It gives
$|\gamma_{j,N}'|\le1/\ell_j$. Define the negative half by conjugate
reflection. Then $B_N$ is continuous on $\T$, with
\begin{equation}\label{eq:BN-modulus}
 B_N(-t)=\overline{B_N(t)},\qquad
 1\le|B_N(t)|\le1+4\delta,\qquad B_N(0)=B_N(1/2)=1.
\end{equation}
Leading phase derivatives match at every join, including the periodic
endpoints. All lie in
$[\delta/(16H),1-\delta/(16H)]$.

On each gap piece, the amplitude $f_N=m_je(\gamma_{j,N})$ satisfies
$\norm{f_N}_\infty+\int|f_N'|\le3+4\pi<16$.
The quadratic estimate gives scaled outer cost below $\delta$, including
reflections. At most one assigned middle derivative interval is within
$\rho$ of $x=k/N$. Its piece and reflection cost at most $192\sqrt{L_g}$.
Every other pair is nonstationary and costs at most
$96/(\rho\sqrt N)$, by one integration by parts with monotone derivative.

Let $C_{M,\Sigma}$ and $C_{NS,\Sigma}$ be the sums of the constants in
Lemma~\ref{lem:stationary} over the main cutoff integrals. Put
\begin{equation}\label{eq:Ncoeff}
 N_{\rm coeff}=\max\left\{1,
 \ceil{\frac{2C_{M,\Sigma}}\delta},
 \ceil{\left(\frac{2C_{NS,\Sigma}}\delta\right)^2},
 \ceil{\left(\frac{96M}{\rho\delta}\right)^2}\right\}.
\end{equation}
For $N\ge N_{\rm coeff}$, coherence, endpoints and gaps give
\begin{equation}\label{eq:interior-coeff}
 \sqrt N\,|\widehat B_N(k)|
 \le1+5\delta+192\sqrt{3.5\delta}<1+500\sqrt\delta,
 \qquad 0\le k<N.
\end{equation}

\subsection{A summable exterior tail}

Partition all main intervals, the three pieces of each gap and their
reflections. On a piece $p$, write $B_N=f_{p,N}e(N\phi_p)$, and let
$F_{p,j}$ bound $\norm{f_{p,N}^{(j)}}_\infty$ for $j\le2$ and
$B_{p,j}$ bound $\norm{\phi_p^{(j)}}_\infty$ for $j=2,3$.
Set $c=\delta/(32H)$. For an exterior integer $k<0$ or $k\ge N$,
$|N\phi_p'-k|\ge c(N+|k|)$. Define
\begin{align}
 A_p&=F_{p,1}/c^2+F_{p,0}B_{p,2}/c^3,\notag\\
 A_p'&=F_{p,2}/c^2+3F_{p,1}B_{p,2}/c^3
          +F_{p,0}B_{p,3}/c^3+3F_{p,0}B_{p,2}^2/c^4,\notag\\
 K_{\rm tail}&=2\sum_p(2A_p+\ell_pA_p').              \label{eq:Ktail}
\end{align}
The first integration-by-parts boundary terms cancel, by continuity of
$B_N$ and matching of the leading derivatives. They also cancel at the
periodic endpoints, since $B_N=1$, $\phi'=1/8$, and $k$ is an integer.
The second integration involves, with $V=N\phi'-k$,
\[
 \frac{f'}{V^2}-\frac{fV'}{V^3},\qquad
 \frac{f''}{V^2}-\frac{3f'V'}{V^3}-\frac{fV''}{V^3}
                        +\frac{3f(V')^2}{V^4}.
\]
Thus \eqref{eq:Ktail} gives
\begin{equation}\label{eq:tail}
 \sum_{k<0\text{ or }k\ge N}|\widehat B_N(k)|\le K_{\rm tail}/N.
\end{equation}
One uses $\sum_{k<0\text{ or }k\ge N}(N+|k|)^{-2}\le2/N$.
The Fourier series is absolutely convergent and equals $B_N$ by
continuity and uniqueness. Conjugate symmetry makes every coefficient real.

\section{An explicit length threshold}\label{sec:threshold}

We now bound every datum used in Section~\ref{sec:analytic} by a fixed
power of one explicit scalar. This also specifies a threshold smaller than
the simpler triple exponential in Theorem~\ref{thm:main}.

For $0<\varepsilon<1$ set
\begin{equation}\label{eq:precision}
 m_\varepsilon=\ceil{2\log_2(2^{20}/\varepsilon)},\qquad
 \delta=2^{-4m_\varepsilon}.
\end{equation}
Using \eqref{eq:aux-data}--\eqref{eq:spreading-T}, define
\begin{align}
 \overline R&=\ceil{K\bigl(T(P_0+1)+1\bigr)^K},&
 \overline m&=\ceil{d+K/2},\notag\\
 \overline U&=\ceil{\max\{H_0,T(P_0+1)/2\}},&
 w_*&=\frac1{\tau T^K(1+\delta)^3},\notag\\
 \overline L&=\ceil{64\overline R^2/\delta},&
 \overline p&=\frac1{2\overline L},\notag\\
 \overline\eta&=\frac\delta{32\overline R},&
 \overline C&=(24\overline R\,\overline L+18\overline R^2)
                      \overline p^{-2\overline R},\notag\\
 \overline A&=\overline m!\,\overline U^{\overline m}.
                                                        \label{eq:master-data}
\end{align}
Put
\begin{align}
 \overline Q=\max\Biggl\{&4,\overline A,
       [2\overline R(\overline C+3\overline R^2+1)]^8,
       [(1+8\overline L)/\overline\eta]^8,\notag\\
       &\left[\frac{2^{20}\overline R}
                         {\overline p\,\overline\eta}\right]
                                      ^{2^{20}\overline R^2}\Biggr\},
                                                        \label{eq:Qbar}\\
 \overline q&=2\ceil{\overline Q},\notag\\
 Z&=2^{64}\bigl(\overline q+\overline R+\delta^{-1}
                        +w_*^{-1}+\overline\eta^{-1}+1\bigr),\notag\\
 N_*(\varepsilon)&=\ceil{Z^{10^7}}.                  \label{eq:Nstar}
\end{align}
Every symbol on the right is an explicit function of $\varepsilon$.
In particular, no instance-specific minimum separation or unspecified
existence threshold occurs in these formulas.

\subsection{All packing choices fit the envelope}

For the actual auxiliary support, write $|A|=2r$. Then
$r\le\overline R$, $m\le\overline m$, and all frequency coordinates
are at most $\overline U$ in absolute value. Every relevant nonzero minor
has absolute value at most $\overline m!\,\overline U^{\overline m}$.
The bounds in Section~\ref{sec:packing} give
$L\le\overline L$, $p\ge\overline p$, and
$\eta\ge\overline\eta$. Therefore every threshold in
\eqref{eq:packing-threshold} is at most $\overline Q$.
Bertrand's postulate gives a prime
\[
 \ceil{\overline Q}<q\le2\ceil{\overline Q}=\overline q.
\]
Proposition~\ref{prop:packing} applies to this prime. All corresponding
centers have denominator $q$, $H\le q$, and \eqref{eq:gap-lengths} holds.
The vector $\theta_h$ can be obtained by solving its finite-field linear
system; this step introduces no larger denominator.

\subsection{Fixed powers control the analytic constants}

From the definition of $Z$ we may use the following loose inequalities:
\begin{equation}\label{eq:basic-Z}
 \begin{gathered}
 H, |A|, q, r, \delta^{-1},\ w_*^{-1},\ \eta^{-1}\le Z,
 \quad M\le Z^2,\quad\rho^{-1}\le Z^3,\\
 \sigma^{-1}\le Z^4,\quad e_0^{-1}\le Z^3,
 \quad(w_*\sigma^2)^{-1}\le Z^9,
 \quad\ell_j^{-1}\le Z^2.
 \end{gathered}
\end{equation}
The sums of the absolute coefficients of $\zeta$ and its first six
derivatives are less than $2^{50}$. This follows directly by expanding
$12012\int_0^s t^6(1-t)^6\,dt$ and differentiating. Thus
\[
 \norm{\chi_h^{(j)}}_\infty\le Z^{3j+1},\qquad
 \norm{G_h^{(j)}}_\infty\le Z^{4j+1}\quad(1\le j\le6).
\]
Leibniz's rule bounds the first six derivatives of $T_{a,h}$ by $Z^{21}$.

Let $\mathcal B_{n,k}$ be the exponential partial Bell polynomial.
Inverse differentiation gives bounds for $X=T_{a,h}^{-1}$ through
\begin{align*}
 J_1&=(w_*\sigma^2)^{-1},\\
 J_n&=(w_*\sigma^2)^{-1}
       \sum_{k=2}^n D_k\mathcal B_{n,k}(J_1,\ldots,J_{n-k+1}),
       \qquad 2\le n\le6,
\end{align*}
where $D_k\le Z^{21}$ bounds $T_{a,h}^{(k)}$.
These imply
\begin{equation}\label{eq:inverse-Z}
 J_n\le Z^{64n^2},\qquad1\le n\le6.
\end{equation}
Indeed a Bell monomial of weight $n$ and $k\ge2$ factors has sum of
squared derivative orders at most $(n-k+1)^2+k-1\le n^2-2$.
The saved $128$ exponent units absorb at most $30$ from the prefactor
and $D_k$, and at most one from the finite coefficient sum. At least
$97$ units of slack remain. The case $n=1$ follows from
\eqref{eq:basic-Z}.

For a main cutoff integral in Lemma~\ref{lem:stationary}, we can take
\begin{equation}\label{eq:stationary-Z}
 b=1,\qquad B=Z^{2500},\qquad
 F_j=Z^{1500}\ (0\le j\le4),\qquad h^{-1}\le Z^{20}.
\end{equation}
Here $w_a<1$ and $\chi_h\le1$ give $|\psi''|\ge1$; the phase
derivative bound follows from $J_6$. For the amplitude
$c_a w_a^{-1/2}e(\sgn(\lambda_a)/8)G_h(X)$, the Bell monomials have
weight at most four, giving an exponent below $64\cdot4^2+17+2<1500$.
The support is at distance at least $w_*\sigma^2e_0/2$ from each main
endpoint, giving the stated $h$ bound.

Substitution into the finite recurrences of Lemma~\ref{lem:stationary}
gives the following conservative bookkeeping bounds:
\begin{align*}
 a_k&\le Z^{100000} &(0\le k\le5),\qquad
 Y_k&\le Z^{100002} &(1\le k\le5),\\
 R_k&\le Z^{500000} &(0\le k\le4),\qquad
 V_{0,k}&\le Z^{501501},\\
 V_{a,k}&\le Z^{501501+500001a} &(a\le4),\qquad
 C_M&\le Z^{3000000},\quad C_{NS}\le Z^{5000}.
\end{align*}
For example the five $a_k$ steps at worst double the preceding exponent
and add a constant, starting at $2500$. Each $R_k$ step adds at most
$2500+100002+1$. The $V$ bounds follow by direct multiplication and
one extra power for each finite sum. There are at most $|A|H\le Z^2$
main integrals. Formula~\eqref{eq:Ncoeff} therefore yields
\begin{equation}\label{eq:Ncoeff-Z}
 N_{\rm coeff}\le Z^{4000000}.
\end{equation}

For gap pieces, \eqref{eq:avoidance} bounds each inverse endpoint
difference by $768M^2\le Z^5$. Formula~\eqref{eq:outer-length} then
bounds the inverse outer lengths by $Z^{13}$. All gap curvatures are
at most $Z^{13}$; for a middle piece use its length at least
$3\ell_j/4$. The gap amplitudes have the uniform derivative bounds
\[
 F_0=2,\qquad F_1\le17/\ell_j,\qquad F_2\le144/\ell_j^2,
\]
which are at most $Z,Z^3,Z^5$. Main amplitudes are constant, and their
phase derivatives through order three are at most $Z^{2500}$.
Since $c^{-1}\le Z^3$ and there are at most $Z^3$ pieces, direct
substitution in \eqref{eq:Ktail} gives
\begin{equation}\label{eq:Ktail-Z}
 K_{\rm tail}\le Z^{6000},\qquad K_{\rm tail}/\delta\le Z^{6001}.
\end{equation}
Thus both analytic thresholds are bounded by $N_*(\varepsilon)$.

\subsection{Projection and rounding to real signs}

For $N\ge N_*(\varepsilon)$ put
\[
 S_\delta=1+500\sqrt\delta,\qquad
 Y_k=\frac{\sqrt N\,\widehat B_N(k)}{S_\delta}\quad(0\le k<N),
 \qquad U_Y(t)=\frac1{\sqrt N}\sum_{k=0}^{N-1}Y_ke(kt).
\]
By \eqref{eq:interior-coeff} the $Y_k$ are real and lie in $[-1,1]$.
Equations~\eqref{eq:tail} and \eqref{eq:Ktail-Z} give
$\norm{U_Y-B_N/S_\delta}_\infty\le\delta$.
Consequently Parseval implies
\[
 \frac1N\sum_kY_k^2\ge(S_\delta^{-1}-\delta)^2,
 \qquad
 \frac\mu N=\frac1{2N}\sum_k(1-|Y_k|)
 \le\frac{1-S_\delta^{-2}}2+\delta
 \le501\sqrt\delta=:q_\delta<\frac12.
\]
Lemma~\ref{lem:rounding}, monotonicity of $q\log(80/q)$ on $(0,1/2]$,
and $\log(80/q_\delta)<2m_\varepsilon$ give real signs with
\begin{align*}
 \norm{U_\epsilon-U_Y}_\infty
 &\le C_R\left(N^{-1/2}+\sqrt{q_\delta\log(80/q_\delta)}\right)\\
 &\le C_RN^{-1/2}+32C_Rm_\varepsilon2^{-m_\varepsilon},
 \qquad U_\epsilon(t)=\frac1{\sqrt N}\sum_{k=0}^{N-1}\epsilon_ke(kt).
\end{align*}
Set $x=2^{20}/\varepsilon$. Then $m_\varepsilon\le x$ and
$2^{-m_\varepsilon}\le x^{-2}$, so
\[
 32C_Rm_\varepsilon2^{-m_\varepsilon}\le\varepsilon/32.
\]
Each of $500\cdot2^{-2m_\varepsilon}$ and
$5\cdot2^{-4m_\varepsilon}$ is also at most $\varepsilon/32$.
Moreover $(2C_R/\varepsilon)^2\le\delta^{-1}\le Z\le N$.
Using \eqref{eq:BN-modulus}, the lower modulus deficit is at most
$500\sqrt\delta+\delta$ plus the rounding error, and the upper excess
is at most $5\delta$ plus that error. Both are therefore at most
\[
 500\cdot2^{-2m_\varepsilon}+5\cdot2^{-4m_\varepsilon}
 +32C_Rm_\varepsilon2^{-m_\varepsilon}+C_R/\sqrt N
 \le\varepsilon.
\]
This proves the desired two-sided estimate whenever $N\ge N_*(\varepsilon)$.
All parameters other than the harmless affine endpoint corrections were
fixed before $N$ was chosen. There is no divisibility restriction on $N$.

\section{The triple exponential bound}\label{sec:tower}

It remains to bound \eqref{eq:Nstar} by the threshold in
Theorem~\ref{thm:main}. Put $t=\log(1/\delta)$. The precision choice
\eqref{eq:precision} ensures $t\ge160\log2>110$ for the entire range
$0<\varepsilon<1$. Directly from the auxiliary formulas,
\begin{align*}
 d&\le16t,&\log K&\le110t^2,&\log P_0&\le150t^2,\\
 \log G_0&\le250t^2,&\log\tau^{-1}&\le130t^2,
 &\log\alpha^{-1}&\le170t^2,\\
 \log R_0&\le3000t^2,&\log\beta_*^{-1}&\le e^{130t^2},
 &\log T&\le e^{200t^2}.
\end{align*}
Here are details of the larger terms. The definitions give
$\log x_0^{-1}\le171t^2$, $\log q_0^{-1}\le684t^2$, and
$\log h_0^{-1}\le685t^2$. Thus
$\log M_2,\log E_1\le2056t^2$, which bounds the maximum defining
$R_0$ as displayed. In
\[
 \log\beta_*^{-1}
 =\log\tau^{-1}+2k_0\log2+\log G_0+(K-1)\log P_0,
\]
the last term is at most $150t^2e^{110t^2}$ and absorbs the others into
$e^{130t^2}$. Taking the ceiling in $T$ is absorbed by $e^{200t^2}$.

Expanding \eqref{eq:master-data} now yields
\begin{align*}
 \log\max(2,\overline R,\overline m,\overline U,w_*^{-1},\overline A)
       &\le e^{500t^2},\\
 \log\max(2,\overline L,\overline p^{-1},\overline\eta^{-1})
       &\le e^{510t^2}.
\end{align*}
For example $\log(T(P_0+1)+1)\le e^{201t^2}$, so
$\log\overline R\le e^{312t^2}$. Also $\overline m\le K$, and
$\log(\overline m!)\le\overline m\log\overline m$ handles the
factorial in $\overline A$.

The logarithms of $\overline C$ and of the matching term in
\eqref{eq:Qbar} are bounded by $\exp(\exp(600t^2))$. Their potentially
largest contributions are respectively
\[
 2\overline R\log\overline p^{-1},\qquad
 2^{20}\overline R^2
       \log\left(\frac{2^{20}\overline R}
                             {\overline p\,\overline\eta}\right).
\]
The remaining finite sums, powers and ceilings in
\eqref{eq:Qbar}--\eqref{eq:Nstar} fit, with room, within
\begin{equation}\label{eq:tower-delta}
 N_*(\varepsilon)\le\Ethree(1000t^2).
\end{equation}
These estimates use $t>110$ throughout, so no unspecified small-error
cutoff is needed. Finally,
\[
 m_\varepsilon\le41+2\log_2(1/\varepsilon),\qquad
 t\le164\log2+8\log(1/\varepsilon)
       \le172\log(2/\varepsilon).
\]
Since $1000\cdot172^2=29{,}584{,}000<10^{12}$,
\eqref{eq:tower-delta} proves Theorem~\ref{thm:main}.

\section*{Scope}

The matching estimate is proved here, including the concentration
inequality it needs. All gap separations and analytic thresholds are given
by explicit numerical constants. The external results used are the
Lovett--Meka partial-coloring theorem and Bertrand's postulate, both cited
rather than re-proved. The triple exponential bounds the length
that suffices; it says nothing about the running time of any algorithm.
No polynomial at these lengths has been computed. The signs may differ from
one length to the next; the paper does not construct a single infinite
sequence of signs.

\section*{AI assistance}

The arguments in this paper were produced with the assistance of AI models
and checked by separate AI proof audits. Agreement between AI audits is not
independent human verification, and the paper has not been peer reviewed.

\begin{thebibliography}{99}

\bibitem{alonkimspencer}
N. Alon, J.-H. Kim and J. Spencer, \emph{Nearly perfect matchings in
regular simple hypergraphs}, Israel Journal of Mathematics \textbf{100}
(1997), 171--187.

\bibitem{alonyuster}
N. Alon and R. Yuster, \emph{On a hypergraph matching problem}, Graphs and
Combinatorics \textbf{21} (2005), no.~4, 377--384.
\url{https://doi.org/10.1007/s00373-005-0628-x}

\bibitem{kkmo}
D. Y. Kang, D. K\"uhn, A. Methuku and D. Osthus, \emph{New bounds on the
size of nearly perfect matchings in almost regular hypergraphs}, Journal of
the London Mathematical Society (2) \textbf{108} (2023), no.~5, 1701--1746.
\url{https://arxiv.org/abs/2010.04183}

\bibitem{bbmst}
P. Balister, B. Bollob\'as, R. Morris, J. Sahasrabudhe and M. Tiba,
\emph{Flat Littlewood polynomials exist}, Annals of Mathematics
\textbf{192} (2020), no.~3, 977--1004.
\url{https://doi.org/10.4007/annals.2020.192.3.6}

\bibitem{bombieribourgain}
E. Bombieri and J. Bourgain, \emph{On Kahane's ultraflat polynomials},
Journal of the European Mathematical Society \textbf{11} (2009), no.~3,
627--703. \url{https://doi.org/10.4171/JEMS/163}

\bibitem{erdelyi}
T. Erd\'elyi, \emph{On an Erd\H{o}s problem about the maximum modulus of
Littlewood polynomials on the unit circle}, arXiv:2608.00744v1, 2026.
\url{https://arxiv.org/abs/2608.00744v1}

\bibitem{kahane}
J.-P. Kahane, \emph{Sur les polyn\^omes \`a coefficients unimodulaires},
Bulletin of the London Mathematical Society \textbf{12} (1980), no.~5,
321--342. \url{https://doi.org/10.1112/blms/12.5.321}

\bibitem{pippengerspencer}
N. Pippenger and J. Spencer, \emph{Asymptotic behavior of the chromatic
index for hypergraphs}, Journal of Combinatorial Theory, Series A
\textbf{51} (1989), no.~1, 24--42.
\url{https://doi.org/10.1016/0097-3165(89)90074-5}

\bibitem{main}
OpenAI, \emph{Ultraflat real Littlewood polynomials}, preprint,
October 5, 2026.
\href{https://github.com/openai/math/tree/main/preprints/Ultraflat-real-Littlewood-polynomials-October-5-2026}{Public source repository}.

\bibitem{predecessor}
OpenAI, \emph{Nearly minimal maxima and positive minima of Littlewood
polynomials}, preprint, October 5, 2026.
\href{https://github.com/openai/math/tree/main/preprints/Nearly-minimal-maxima-and-positive-minima-of-Littlewood-polynomials-October-5-2026}{Public source repository}.

\bibitem{lovettmeka}
Shachar Lovett and Raghu Meka, \emph{Constructive discrepancy minimization
by walking on the edges}, arXiv:1203.5747v2, Theorem 4.
\href{https://arxiv.org/pdf/1203.5747v2}{Versioned theorem source}.

\end{thebibliography}
\end{document}
